% 康奈尔听课笔记模板（配色取自王道课件）
\documentclass[10.5pt,a4paper,oneside]{article}
\usepackage[a4paper,top=2.3cm,bottom=2.2cm,left=5.7cm,right=1.5cm,
  marginparwidth=4.25cm,marginparsep=0.55cm,headheight=14pt,footskip=1.1cm]{geometry}
\usepackage{xeCJK}
\usepackage{xeCJKfntef}
\xeCJKDeclareCharClass{FullLeft}{"201C,"2018}
\xeCJKDeclareCharClass{FullRight}{"201D,"2019}
\xeCJKDeclareCharClass{CJK}{"2460->"24FF,"2605,"2606}
\setmainfont{Noto Sans CJK SC}[BoldFont={Noto Sans CJK SC Bold},ItalicFont={Noto Sans CJK SC},ItalicFeatures={FakeSlant=0.18}]
\setsansfont{Noto Sans CJK SC}[BoldFont={Noto Sans CJK SC Bold}]
\setmonofont{DejaVu Sans Mono}[Scale=0.9]
\setCJKmainfont{Noto Sans CJK SC}[BoldFont={Noto Sans CJK SC Bold},ItalicFont={Noto Sans CJK SC},ItalicFeatures={FakeSlant=0.18}]
\setCJKsansfont{Noto Sans CJK SC}[BoldFont={Noto Sans CJK SC Bold}]
\setCJKmonofont{Noto Sans Mono CJK SC}[Scale=0.9]
\linespread{1.32}
\usepackage{amsmath,amssymb}
\usepackage[dvipsnames,table]{xcolor}
\usepackage{graphicx}
\usepackage{adjustbox}
\usepackage{tikz}
\usetikzlibrary{calc}
\usepackage[most]{tcolorbox}
\usepackage{eso-pic}
\usepackage{fancyhdr}
\usepackage{titlesec}
\usepackage{needspace}
\usepackage{enumitem}
\usepackage{xltabular,booktabs,array,colortbl}
\usepackage{marginnote}
\reversemarginpar
\renewcommand*{\raggedleftmarginnote}{}
\renewcommand*{\raggedrightmarginnote}{}
\usepackage{caption}
\usepackage{titletoc}
\usepackage{fancyvrb}
\usepackage{framed}
\usepackage[section]{placeins}
\usepackage{hyperref}
\hypersetup{colorlinks=true,linkcolor=wdblue,urlcolor=wdblue,pdftitle={3.3.3
栈在递归中的应用},pdfauthor={听课笔记}}

% ---------- 课件配色 ----------
\definecolor{wdblue}{HTML}{4472C4}
\definecolor{wdlightblue}{HTML}{5B9BD5}
\definecolor{wdtitle}{HTML}{4A86C8}
\definecolor{wddark}{HTML}{1F3864}
\definecolor{wdorange}{HTML}{ED7D31}
\definecolor{wdgray}{HTML}{A5A5A5}
\definecolor{wdgold}{HTML}{FFC000}
\definecolor{wdgreen}{HTML}{70AD47}
\definecolor{wddeepgreen}{HTML}{3F7F2A}
\definecolor{wdred}{HTML}{C00000}
\definecolor{wdpurple}{HTML}{7030A0}
\definecolor{cuebg}{HTML}{F2F6FC}
\definecolor{codebg}{HTML}{F7F9FC}
\color[HTML]{262626}

% ---------- 行内样式 ----------
\newcommand{\hlred}[1]{\textcolor{wdred}{\bfseries #1}}
\newcommand{\hlblue}[1]{\textcolor{wdblue}{\bfseries #1}}
\newcommand{\hlgreen}[1]{\textcolor{wddeepgreen}{\bfseries #1}}
\newcommand{\hlorange}[1]{\textcolor{wdorange}{\bfseries #1}}
\newcommand{\hlpurple}[1]{\textcolor{wdpurple}{\bfseries #1}}
\newcommand{\uln}[1]{\CJKunderline[format=\color{wdred},thickness=0.9pt,depth=0.28em]{#1}}
\newcommand{\wavy}[1]{\CJKunderwave[format=\color{wdorange}]{#1}}
\newcommand{\hlmark}[1]{\CJKunderline[format=\color{wdgold!55},thickness=0.55em,depth=-0.05em]{#1}}
\newcommand{\tagbox}[2]{\tcbox[on line,colback=#1,colframe=#1,coltext=white,boxsep=0.5pt,left=3pt,right=3pt,top=0.5pt,bottom=0.5pt,arc=2pt,fontupper=\small\bfseries]{#2}}

% ---------- 康奈尔版式背景：左侧线索栏 ----------
\newlength{\cuestripw}\setlength{\cuestripw}{5.15cm}
\AddToShipoutPictureBG{%
  \begin{tikzpicture}[remember picture,overlay]
    \fill[cuebg] (current page.north west) rectangle ($(current page.south west)+(\cuestripw,0)$);
    \draw[wdblue,line width=1.2pt] ($(current page.north west)+(\cuestripw,-1.55cm)$) -- ($(current page.south west)+(\cuestripw,1.55cm)$);
    \node[anchor=west,text=wdblue,font=\small\bfseries] at ($(current page.north west)+(0.75cm,-1.25cm)$) {线索 · 问题};
    \node[anchor=west,text=wdblue,font=\small\bfseries] at ($(current page.north west)+(\cuestripw+0.55cm,-1.25cm)$) {笔记 · 老师主线 + 补充讲解};
  \end{tikzpicture}}

% ---------- 页眉页脚 ----------
\pagestyle{fancy}\fancyhf{}
\renewcommand{\headrulewidth}{0pt}
\fancyhead[R]{\small\color{wdgray}第3章 栈、队列和数组}
\fancyfoot[C]{\small\color{wdgray}\thepage}
\fancyfoot[R]{\small\color{wdorange}3.3.3 栈在递归中的应用 · 听课笔记}

\fancypagestyle{plain}{\fancyhf{}\fancyfoot[C]{\small\color{wdgray}\thepage}\fancyfoot[R]{\small\color{wdorange}3.3.3
栈在递归中的应用 · 听课笔记}}

% ---------- 标题：仿课件的蓝色标题 + 六色条 ----------
\newcommand{\sixbar}{\par\nointerlineskip\vspace{3pt}\centering
  \begin{tikzpicture}\foreach \c [count=\i] in {wdlightblue,wdorange,wdgray,wdgold,wdblue,wdgreen}
     {\fill[\c] ({(\i-1)*1.05},0) rectangle ({(\i-1)*1.05+0.92},0.11);}\end{tikzpicture}\par}
\titleformat{\section}[block]{\centering\Large\bfseries\color{wdtitle}}{}{0pt}{}[\sixbar]
\titlespacing*{\section}{0pt}{16pt}{10pt}
\pretocmd{\section}{\Needspace{8\baselineskip}}{}{}
\titleformat{\subsection}{\Needspace{7\baselineskip}\large\bfseries\color{wddark}}{\textcolor{wdblue}{\rule[-0.15em]{4pt}{1.05em}}\hspace{6pt}}{0pt}{}
\titlespacing*{\subsection}{0pt}{12pt}{5pt}
\titleformat{\subsubsection}{\Needspace{5\baselineskip}\normalsize\bfseries\color{wdblue}}{\textcolor{wdorange}{\large\ensuremath{\blacktriangleright}}\hspace{4pt}}{0pt}{}
\titlespacing*{\subsubsection}{0pt}{8pt}{3pt}
\setcounter{secnumdepth}{0}
\setlength{\parindent}{0pt}\setlength{\parskip}{5pt}
\setlist{nosep,leftmargin=1.5em,itemsep=2pt}
\setlist[itemize,1]{label=\textcolor{wdblue}{\small\ensuremath{\bullet}}}
\setlist[itemize,2]{label=\textcolor{wdorange}{\scriptsize\ensuremath{\blacktriangleright}}}
\setlist[enumerate,1]{label=\textcolor{wdblue}{\bfseries\arabic*.}}
\providecommand{\tightlist}{\setlength{\itemsep}{2pt}\setlength{\parskip}{0pt}}

% ---------- 目录 ----------
\setcounter{tocdepth}{3}
\titlecontents{section}[0pt]{\addvspace{5pt}\bfseries\color{wddark}}{}{}{\titlerule*[0.7pc]{\color{wdgray!60}.}\contentspage}
\titlecontents{subsection}[1.6em]{\addvspace{1pt}\small\color{black!75}}{}{}{\titlerule*[0.7pc]{\color{wdgray!40}.}\contentspage}
\newcommand{\tocback}{\par\vspace{-3pt}{\hfill\footnotesize\hyperlink{toc}{\textcolor{wdgray}{\ensuremath{\uparrow}\,回到目录}}}\par\vspace{-4pt}}

% ---------- 线索栏 ----------
\newcommand{\cueblock}[2]{\Needspace{4\baselineskip}\leavevmode#1{#2}\par\vspace{-\baselineskip}\vspace{-\parskip}}
% 线索直接贴在段首行的左侧（不用 marginnote，避免分页时偶尔落进正文栏）
\newcommand{\cuebox}[1]{\leavevmode\llap{\smash{\begin{minipage}[t]{\marginparwidth}\raggedright\small\linespread{1.18}\selectfont #1\end{minipage}}\hspace{\marginparsep}}}
\newcommand{\cuenote}[1]{\cuebox{{\color{wdblue}\bfseries Q\,}\color{wddark}#1}}
\newcommand{\cuekey}[1]{\cuebox{\color{wdred}\bfseries #1}}

% ---------- 讲解框 ----------
\tcbset{wdbox/.style={enhanced,breakable,arc=2pt,boxrule=0pt,leftrule=3pt,left=7pt,right=7pt,top=5pt,bottom=5pt,
   before skip=7pt,after skip=7pt,fonttitle=\bfseries\small,attach title to upper={\quad},coltitle=#1,
   colframe=#1,colback=#1!6}}
\newtcolorbox{defbox}[1][]{enhanced,breakable,arc=3pt,boxrule=0.8pt,colframe=wdblue,colback=wdblue!4,
   left=8pt,right=8pt,top=6pt,bottom=6pt,fonttitle=\bfseries\small,coltitle=white,colbacktitle=wdblue,
   title={定义#1},before skip=8pt,after skip=8pt}
\newtcolorbox{thmbox}[1]{enhanced,breakable,arc=3pt,boxrule=0.8pt,colframe=wdblue,colback=wdblue!4,
   left=8pt,right=8pt,top=6pt,bottom=6pt,fonttitle=\bfseries\small,coltitle=white,colbacktitle=wdblue,
   title={#1},before skip=8pt,after skip=8pt}
\newtcolorbox{bridgebox}{wdbox=wdblue,title={承上启下}}
\newtcolorbox{plainbox}{wdbox=wddeepgreen,colframe=wdgreen,colback=wdgreen!8,title={大白话}}
\newtcolorbox{deepbox}{wdbox=wdpurple,title={背后的思想}}
\newtcolorbox{exambox}{wdbox=wdred,title={考点}}
\newtcolorbox{pitfallbox}{wdbox=wdorange,title={易错提醒}}
\newtcolorbox{bookbox}{wdbox=wddark,colback=wdgray!10,colframe=wdgray,title={王道书补充}}
\newtcolorbox{teacherbox}{wdbox=wdorange,colback=wdgold!8,colframe=wdgold,coltitle=wdorange,title={老师原话}}
\newtcolorbox{fixbox}{wdbox=wdred,colback=wdred!4,title={勘误 / 说明}}
% 习题笔记用
\newtcolorbox{questionbox}[1][]{enhanced,breakable,arc=3pt,boxrule=0.8pt,colframe=wddark,colback=white,
   left=8pt,right=8pt,top=6pt,bottom=6pt,fonttitle=\bfseries\small,coltitle=white,colbacktitle=wddark,
   title={#1},before skip=10pt,after skip=6pt}
\newtcolorbox{videobox}{wdbox=wdorange,colback=wdgold!8,colframe=wdgold,coltitle=wdorange,title={视频讲解}}
\newtcolorbox{jiexibox}{wdbox=wddark,colback=wdgray!10,colframe=wdgray,title={书上解析}}
\newtcolorbox{addbox}{wdbox=wdpurple,title={我的补充}}
\newtcolorbox{recallbox}{wdbox=wdblue,title={回顾正课}}
\newtcolorbox{watchbox}{enhanced,breakable,arc=3pt,boxrule=1pt,colframe=wdorange,colback=wdgold!12,
   fonttitle=\bfseries\small,coltitle=white,colbacktitle=wdorange,title={建议去看视频},left=8pt,right=8pt,before skip=8pt,after skip=8pt}
\newtcolorbox{summarybox}{enhanced,arc=3pt,
   boxrule=1pt,colframe=wdblue,colback=white,colbacktitle=wdblue,coltitle=white,fonttitle=\bfseries,
   title={总结区 · 用自己的话复述这一节},left=10pt,right=10pt,top=8pt,bottom=8pt,before skip=14pt}

% ---------- 代码 ----------
\usepackage{color}
\usepackage{fancyvrb}
\newcommand{\VerbBar}{|}
\newcommand{\VERB}{\Verb[commandchars=\\\{\}]}
\DefineVerbatimEnvironment{Highlighting}{Verbatim}{commandchars=\\\{\}}
% Add ',fontsize=\small' for more characters per line
\newenvironment{Shaded}{}{}
\newcommand{\AlertTok}[1]{\textcolor[rgb]{1.00,0.00,0.00}{#1}}
\newcommand{\AnnotationTok}[1]{\textcolor[rgb]{0.38,0.63,0.69}{#1}}
\newcommand{\AttributeTok}[1]{\textcolor[rgb]{0.49,0.56,0.16}{#1}}
\newcommand{\BaseNTok}[1]{\textcolor[rgb]{0.77,0.35,0.07}{#1}}
\newcommand{\BuiltInTok}[1]{\textcolor[rgb]{0.00,0.50,0.00}{#1}}
\newcommand{\CharTok}[1]{\textcolor[rgb]{0.63,0.19,0.43}{#1}}
\newcommand{\CommentTok}[1]{\textcolor[rgb]{0.18,0.55,0.23}{#1}}
\newcommand{\CommentVarTok}[1]{\textcolor[rgb]{0.18,0.55,0.23}{#1}}
\newcommand{\ConstantTok}[1]{\textcolor[rgb]{0.77,0.35,0.07}{#1}}
\newcommand{\ControlFlowTok}[1]{\textcolor[rgb]{0.12,0.31,0.71}{#1}}
\newcommand{\DataTypeTok}[1]{\textcolor[rgb]{0.12,0.31,0.71}{#1}}
\newcommand{\DecValTok}[1]{\textcolor[rgb]{0.77,0.35,0.07}{#1}}
\newcommand{\DocumentationTok}[1]{\textcolor[rgb]{0.18,0.55,0.23}{#1}}
\newcommand{\ErrorTok}[1]{\textcolor[rgb]{1.00,0.00,0.00}{#1}}
\newcommand{\ExtensionTok}[1]{#1}
\newcommand{\FloatTok}[1]{\textcolor[rgb]{0.77,0.35,0.07}{#1}}
\newcommand{\FunctionTok}[1]{\textcolor[rgb]{0.44,0.19,0.63}{#1}}
\newcommand{\ImportTok}[1]{\textcolor[rgb]{0.44,0.19,0.63}{#1}}
\newcommand{\InformationTok}[1]{\textcolor[rgb]{0.38,0.63,0.69}{#1}}
\newcommand{\KeywordTok}[1]{\textcolor[rgb]{0.12,0.31,0.71}{#1}}
\newcommand{\NormalTok}[1]{#1}
\newcommand{\OperatorTok}[1]{\textcolor[rgb]{0.20,0.20,0.20}{#1}}
\newcommand{\OtherTok}[1]{\textcolor[rgb]{0.00,0.44,0.13}{#1}}
\newcommand{\PreprocessorTok}[1]{\textcolor[rgb]{0.44,0.19,0.63}{#1}}
\newcommand{\RegionMarkerTok}[1]{#1}
\newcommand{\SpecialCharTok}[1]{\textcolor[rgb]{0.63,0.19,0.43}{#1}}
\newcommand{\SpecialStringTok}[1]{\textcolor[rgb]{0.63,0.19,0.43}{#1}}
\newcommand{\StringTok}[1]{\textcolor[rgb]{0.63,0.19,0.43}{#1}}
\newcommand{\VariableTok}[1]{\textcolor[rgb]{0.10,0.09,0.49}{#1}}
\newcommand{\VerbatimStringTok}[1]{\textcolor[rgb]{0.63,0.19,0.43}{#1}}
\newcommand{\WarningTok}[1]{\textcolor[rgb]{0.38,0.63,0.69}{#1}}
\makeatletter
\@ifundefined{Shaded}{}{\renewenvironment{Shaded}{\def\FrameCommand{\fboxsep=5pt\fboxrule=0.5pt\fcolorbox{wdblue!35}{codebg}}%
   \MakeFramed{\advance\hsize-\width\FrameRestore}\small}{\endMakeFramed}}
\makeatother
\fvset{fontsize=\small}

% ---------- 表格 ----------
\renewcommand{\arraystretch}{1.42}
\newcolumntype{Y}{>{\raggedright\arraybackslash}X}

% ---------- 图片 ----------
% 课件截图做成浮动体：LaTeX 会自动用正文填满页面，避免大块空白
\setcounter{topnumber}{3}\setcounter{bottomnumber}{2}\setcounter{totalnumber}{4}
\renewcommand{\topfraction}{0.85}\renewcommand{\bottomfraction}{0.6}\renewcommand{\textfraction}{0.12}\renewcommand{\floatpagefraction}{0.75}
\setlength{\textfloatsep}{10pt plus 3pt minus 2pt}\setlength{\floatsep}{8pt plus 2pt}
\newcommand{\slidefig}[3]{\begin{figure}[!ht]\centering\begin{tcolorbox}[enhanced,width=#2\linewidth,arc=2pt,boxrule=0.6pt,colframe=wdgray!70,
   colback=white,left=2pt,right=2pt,top=2pt,bottom=2pt,drop fuzzy shadow=black!25]\includegraphics[width=\linewidth]{#1}\end{tcolorbox}
   \vspace{-3pt}{\small\color{wdgray}#3}\end{figure}}
\newcommand{\plainfig}[3]{\par\noindent\begin{minipage}{\linewidth}\centering\includegraphics[width=#2\linewidth]{#1}\par\vspace{2pt}{\small\color{wdgray}#3}\end{minipage}\par\vspace{4pt}}

\begin{document}
% ---------- 封面条（仿课件封面） ----------
\thispagestyle{plain}
\begin{tcolorbox}[enhanced,enlarge left by=-4.8cm,width=\linewidth+4.8cm,arc=0pt,boxrule=0pt,
   colback=wddark,interior style={left color=wddark,right color=wdblue},left=18pt,right=18pt,top=14pt,bottom=12pt,
   overlay={\draw[white,line width=2.5pt] ([xshift=10pt,yshift=-10pt]frame.north west) -- ++(0,-0.9cm);
            \draw[white,line width=2.5pt] ([xshift=10pt,yshift=-10pt]frame.north west) -- ++(1.6cm,0);
            \draw[white,line width=2.5pt] ([xshift=-10pt,yshift=10pt]frame.south east) -- ++(0,0.9cm);
            \draw[white,line width=2.5pt] ([xshift=-10pt,yshift=10pt]frame.south east) -- ++(-1.6cm,0);}]
  {\color{white!85!wdlightblue}\small 数据结构 · 第3章
栈、队列和数组 · 康奈尔听课笔记}\par\vspace{6pt}
  {\color{white}\fontsize{22}{28}\selectfont\bfseries 3.3.3
栈在递归中的应用}\par\vspace{6pt}
  {\color{white!80!wdlightblue}\small 课件 3.3.3 · 正课字幕 P30 · 王道书
§3.3.3（第 92--93 页）}
\end{tcolorbox}
\vspace{4pt}
\hypertarget{ux8fd9ux4e00ux8282ux5728ux8bb2ux4ec0ux4e48}{%
\section{这一节在讲什么}\label{ux8fd9ux4e00ux8282ux5728ux8bb2ux4ec0ux4e48}}

\Needspace{4\baselineskip}

\cuenote{递归和栈有什么关系？}递归算法就是递归函数调用它自身的过程，所以\uln{递归的背后就是函数调用}。这一节先看普通的函数调用背后发生了什么（要用一个栈），再看递归调用时这个栈怎么变化，以及递归的缺点。老师说：这一节不是讲怎么设计递归算法，而是理解\uln{递归算法和栈有什么本质联系}，大致了解函数调用栈的原理就可以。

\hypertarget{ux4e00ux51fdux6570ux8c03ux7528ux80ccux540eux7684ux8fc7ux7a0b}{%
\section{一、函数调用背后的过程}\label{ux4e00ux51fdux6570ux8c03ux7528ux80ccux540eux7684ux8fc7ux7a0b}}

\Needspace{4\baselineskip}

\cuenote{函数调用有什么特点？函数调用栈里存哪三样东西？}\texttt{main}
调用 \texttt{func1}，\texttt{func1} 又调用
\texttt{func2}；\texttt{func2} 执行结束后返回 \texttt{func1}
接着往后执行，\texttt{func1} 结束后再返回
\texttt{main}。函数调用的特点：\hlred{最后被调用的函数最先执行结束}（LIFO），和栈的后进先出异曲同工。

事实上，任何程序运行之前，系统都会开辟一个\textbf{函数调用栈}，保存各个函数在调用过程中必须保存的信息：

\begin{enumerate}
\def\labelenumi{\arabic{enumi}.}
\tightlist
\item
  \textbf{调用返回地址}：被调函数执行结束后应该接着执行哪一句代码，把那句代码的存储地址压入栈中（图中的
  \texttt{\#1}、\texttt{\#2}）；
\item
  \textbf{实参}：如 \texttt{func1} 的参数 \texttt{a}、\texttt{b}；
\item
  \textbf{局部变量}：如 \texttt{func1} 里定义的 \texttt{x}。
\end{enumerate}

\slidefig{assets/3.3.3-p02.jpg}{0.6}{课件：函数调用背后的过程，函数调用栈}

老师补充的几点：

\begin{itemize}
\tightlist
\item
  这个栈其实就是\uln{内存里的某一片存储空间}。被调函数执行结束，就把和它相关的信息弹出栈，也就是释放这片空间，再根据栈顶记录的返回地址，接着执行下一句；
\item
  实参在栈里是另一份数据，所以在 \texttt{func1} 里修改
  \texttt{a}、\texttt{b} 的值，影响不到 \texttt{main} 里的
  \texttt{a}、\texttt{b}；
\item
  编译器会在 \texttt{main} 之前加一些别的代码，所以栈底在 \texttt{main}
  下面还有一些信息（图中的“\ldots\ldots{}”），普通程序员不用关心；
\item
  用 IDE（老师用的是 CLion，VS
  等都可以）在某句代码下断点调试，可以看到此时的函数调用栈：栈顶是
  \texttt{func2}，里面有参数 \texttt{x} 和局部变量
  \texttt{m}、\texttt{n}；点栈里的 \texttt{func1}，能看到
  \texttt{a}、\texttt{b}、\texttt{x}。IDE 里不显示返回地址。
\end{itemize}

\slidefig{assets/3.3.3-p03.jpg}{0.6}{课件：在 IDE 中调试，观察函数调用栈}

\hypertarget{ux4e8cux6808ux5728ux9012ux5f52ux4e2dux7684ux5e94ux7528}{%
\section{二、栈在递归中的应用}\label{ux4e8cux6808ux5728ux9012ux5f52ux4e2dux7684ux5e94ux7528}}

\Needspace{4\baselineskip}

\cuenote{什么问题适合用递归？实现递归需要哪两样东西？}适合用“递归”算法解决的问题：\hlred{可以把原始问题转换为属性相同，但规模较小的问题}。

\begin{itemize}
\tightlist
\item
  \textbf{例 1}：计算正整数的阶乘
  \texttt{n!}，\texttt{n!\ =\ n\ ×\ (n-1)!}，问题规模从 n 变成
  n-1，属性相同，都是求阶乘；
\item
  \textbf{例 2}：求斐波那契数列，把规模 n 拆成规模 n-1 和 n-2 两部分。
\end{itemize}

实现一个递归算法需要两个比较重要的东西：\textbf{递归表达式}（递归体）和\textbf{边界条件}（递归出口）。

\slidefig{assets/3.3.3-p04.jpg}{0.55}{课件：递归表达式（递归体）和边界条件（递归出口）}

\hypertarget{ux4f8b-1ux9012ux5f52ux7b97ux6cd5ux6c42ux9636ux4e58}{%
\subsection{例
1：递归算法求阶乘}\label{ux4f8b-1ux9012ux5f52ux7b97ux6cd5ux6c42ux9636ux4e58}}

\begin{Shaded}
\begin{Highlighting}[]
\CommentTok{// 计算正整数 n!}
\DataTypeTok{int}\NormalTok{ factorial}\OperatorTok{(}\DataTypeTok{int}\NormalTok{ n}\OperatorTok{)} \OperatorTok{\{}
    \ControlFlowTok{if} \OperatorTok{(}\NormalTok{n }\OperatorTok{==} \DecValTok{0} \OperatorTok{||}\NormalTok{ n }\OperatorTok{==} \DecValTok{1}\OperatorTok{)}
        \ControlFlowTok{return} \DecValTok{1}\OperatorTok{;}
    \ControlFlowTok{else}
        \ControlFlowTok{return}\NormalTok{ n }\OperatorTok{*}\NormalTok{ factorial}\OperatorTok{(}\NormalTok{n }\OperatorTok{{-}} \DecValTok{1}\OperatorTok{);}    \CommentTok{// 第 187 行}
\OperatorTok{\}}

\DataTypeTok{int}\NormalTok{ main}\OperatorTok{()} \OperatorTok{\{}
    \CommentTok{// ... 其他代码}
    \DataTypeTok{int}\NormalTok{ x }\OperatorTok{=}\NormalTok{ factorial}\OperatorTok{(}\DecValTok{10}\OperatorTok{);}              \CommentTok{// 第 192 行}
\NormalTok{    printf}\OperatorTok{(}\StringTok{"奥利给！"}\OperatorTok{);}
\OperatorTok{\}}
\end{Highlighting}
\end{Shaded}

递归调用时，函数调用栈可称为\hlred{\textbf{递归工作栈}}，本质上和函数调用栈是同一个东西。

\begin{enumerate}
\def\labelenumi{\arabic{enumi}.}
\tightlist
\item
  \texttt{main} 调用 \texttt{factorial(10)}：压入返回地址
  \texttt{\#192}（结束后要把结果赋给 \texttt{x}）和实参
  \texttt{n=10}，这是第 1 层；
\item
  \texttt{n=10} 不等于 0 也不等于 1，执行第 187 行，递归调用
  \texttt{factorial(9)}：压入返回地址 \texttt{\#187}（回来后要把返回值和
  n 相乘）和 \texttt{n=9}，第 2 层；
\item
  依此类推，问题规模逐渐收敛，直到第 10 层
  \texttt{n=1}，满足边界条件，直接 \texttt{return\ 1}；
\item
  最深一层返回后，弹出它的信息，回到上一层（\texttt{n=2}）继续执行第 187
  行，返回 \texttt{2×1}；再逐层返回，直到第 1 层算出
  \texttt{10×9!}，返回给 \texttt{main}，赋给 \texttt{x}。
\end{enumerate}

所以：\uln{每进入一层递归，就将递归调用所需信息压入栈顶；每退出一层递归，就从栈顶弹出相应信息}。

\slidefig{assets/3.3.3-p05.jpg}{0.6}{课件：递归算法求阶乘，递归工作栈}

\textbf{缺点：太多层递归可能会导致栈溢出。} 如果传入 10 万，栈里就要压入
10
万条信息，而内存资源有限，系统开辟的函数调用栈也有存储上限。课件上“再次思考：递归算法的空间复杂度”：这也就是绪论里讲的\uln{递归带来空间复杂度升高}，递归层数越多，空间复杂度越高。

\begin{recallbox}

正课 1.2.3 讲空间复杂度时用 \texttt{loveYou(n)}
讲过：每层调用单独开一块空间存本层的参数、局部变量和返回地址，空间复杂度＝递归调用的深度（每层所需空间为常量时）。这里的递归工作栈就是这部分内存开销的来源。

\end{recallbox}

由于递归算法背后就是用一个栈实现的，\uln{可以自定义栈将递归算法改造成非递归算法}，类似的题目会在课后习题中遇到。

\slidefig{assets/3.3.3-p06.jpg}{0.55}{课件：在 IDE 中观察递归工作栈}

\begin{addbox}

\textbf{factorial(3)
的栈变化}（我自己写的简化版，栈从左往右是栈底到栈顶，省略 \texttt{main}
下面的信息）：

{\rowcolors{2}{wdblue!5}{white}\arrayrulecolor{wdblue!35}
\begin{xltabular}{\linewidth}{>{\hsize=0.720\hsize\raggedright\arraybackslash}X>{\hsize=1.280\hsize\raggedright\arraybackslash}X}
\rowcolor{wdblue}\textcolor{white}{\bfseries 执行到} & \textcolor{white}{\bfseries 递归工作栈} \\
\endfirsthead
main 调用 factorial(3) & main; \#192, n=3 \\ \hline
n=3，调用 factorial(2) & main; \#192, n=3; \#187, n=2 \\ \hline
n=2，调用 factorial(1) & main; \#192, n=3; \#187, n=2; \#187, n=1 \\ \hline
n=1 满足边界条件，return 1，弹出 & main; \#192, n=3; \#187, n=2 \\ \hline
n=2 层算 2×1=2，return 2，弹出 & main; \#192, n=3 \\ \hline
n=3 层算 3×2=6，return 6，弹出 & main（x=6） \\ \hline
\end{xltabular}\par\vspace{4pt}}

递归层数 3，栈里最多同时有 3 层递归的信息；\texttt{factorial(n)} 就是 n
层，空间复杂度 O(n)。

\end{addbox}

\hypertarget{ux4f8b-2ux9012ux5f52ux7b97ux6cd5ux6c42ux6590ux6ce2ux90a3ux5951ux6570ux5217}{%
\subsection{例
2：递归算法求斐波那契数列}\label{ux4f8b-2ux9012ux5f52ux7b97ux6cd5ux6c42ux6590ux6ce2ux90a3ux5951ux6570ux5217}}

\begin{Shaded}
\begin{Highlighting}[]
\DataTypeTok{int}\NormalTok{ Fib}\OperatorTok{(}\DataTypeTok{int}\NormalTok{ n}\OperatorTok{)} \OperatorTok{\{}
    \ControlFlowTok{if} \OperatorTok{(}\NormalTok{n }\OperatorTok{==} \DecValTok{0}\OperatorTok{)}
        \ControlFlowTok{return} \DecValTok{0}\OperatorTok{;}
    \ControlFlowTok{else} \ControlFlowTok{if} \OperatorTok{(}\NormalTok{n }\OperatorTok{==} \DecValTok{1}\OperatorTok{)}
        \ControlFlowTok{return} \DecValTok{1}\OperatorTok{;}
    \ControlFlowTok{else}
        \ControlFlowTok{return}\NormalTok{ Fib}\OperatorTok{(}\NormalTok{n }\OperatorTok{{-}} \DecValTok{1}\OperatorTok{)} \OperatorTok{+}\NormalTok{ Fib}\OperatorTok{(}\NormalTok{n }\OperatorTok{{-}} \DecValTok{2}\OperatorTok{);}
\OperatorTok{\}}

\DataTypeTok{int}\NormalTok{ main}\OperatorTok{()} \OperatorTok{\{}
    \CommentTok{// ... 其他代码}
    \DataTypeTok{int}\NormalTok{ x }\OperatorTok{=}\NormalTok{ Fib}\OperatorTok{(}\DecValTok{4}\OperatorTok{);}
\NormalTok{    printf}\OperatorTok{(}\StringTok{"奥利给！"}\OperatorTok{);}
\OperatorTok{\}}
\end{Highlighting}
\end{Shaded}

\texttt{Fib(4)} 先调用 \texttt{Fib(3)}，进入更深一层；\texttt{Fib(3)}
还没到边界，先调用 \texttt{Fib(2)}；\texttt{Fib(2)} 调用
\texttt{Fib(1)}，n=1 满足边界条件，返回，弹出栈顶；\texttt{Fib(2)}
再调用 \texttt{Fib(0)}，返回；\texttt{Fib(2)} 两个调用都有返回值了，返回
\texttt{Fib(3)}；\texttt{Fib(3)} 再调用 \texttt{Fib(1)}，返回后
\texttt{Fib(3)} 返回
\texttt{Fib(4)}\ldots\ldots 就这样逐层完成递归调用。

\textbf{缺点：可能包含很多重复计算。}
在这个过程中，\uln{\texttt{Fib(2)} 被计算了两次}，\texttt{Fib(1)}、\texttt{Fib(0)}
也被重复了多次。这也是递归算法不太高效的原因之一。

\slidefig{assets/3.3.3-p07.jpg}{0.6}{课件：递归算法求斐波那契数列，可能包含很多重复计算}

\begin{addbox}

数一下 \texttt{Fib(4)} 的递归调用树：\texttt{Fib(2)} 2
次，\texttt{Fib(1)} 3 次，\texttt{Fib(0)} 2 次，一共调用 9
次。栈里最多同时有 4 层（n=4、3、2、1 或 0），即递归深度。

\end{addbox}

\begin{bookbox}

王道书第 92--93 页
§3.3.3：若一个函数、过程或数据结构在其定义中直接或间接地引用了自身，则称其为递归定义，简称递归。递归通过将原问题逐层分解为规模更小但结构相同的子问题来求解；仅需少量代码即可简洁地表达解题过程中反复出现的相同计算模式，显著减少程序代码量。然而，在一般情况下，递归的执行效率并不高。

书上以斐波那契数列为例（代码与课件例 2 相同，书上函数名为
\texttt{F}，注释标出两个 \texttt{return\ 0/1}
是“边界条件”，\texttt{return\ F(n-1)+F(n-2)}
是“递归表达式”）。递归定义不能是循环或无终止的，必须满足两个条件：

\begin{itemize}
\tightlist
\item
  \textbf{递归体}（递归表达式）：描述问题如何分解为更小规模的同类子问题；
\item
  \textbf{边界条件}（递归出口）：确保递归在有限步内终止。
\end{itemize}

递归的精髓在于能否将原问题转化为\uln{属性相同但规模更小}的问题。

考点追踪（栈在函数调用中的作用和工作原理，2015、2017）：在递归调用过程中，系统为每一层调用的\uln{返回地址、局部变量及形参}等信息分配存储空间，这些信息被保存在递归工作栈中。因此，递归深度过大时容易导致\textbf{栈溢出}。此外，递归效率低下的主要原因是\textbf{存在大量重复计算}。以
n=5 为例，见图 3.16：\texttt{F(3)} 被计算 2 次，\texttt{F(2)} 被计算 3
次，\texttt{F(1)} 被调用 5 次，\texttt{F(0)} 被调用 3 次。

递归虽然代码简洁、易于理解，但存在明显的性能缺陷。第 5
章讨论树的遍历时，采用递归可使代码极为简洁。\uln{任何递归算法均可转换为非递归算法}，通常需借助显式栈来模拟递归调用过程。

\end{bookbox}

\slidefig{assets/3.3-book-fig3.16.png}{0.6}{王道书图 3.16：F(5) 的递归执行过程}

\begin{addbox}

我按图 3.16 核对过次数：\texttt{F(5)} 的调用树中 \texttt{F(3)} 2
次、\texttt{F(2)} 3 次、\texttt{F(1)} 5 次、\texttt{F(0)} 3
次，和书上一致（也可以这样记：\texttt{F(1)} 被调用的次数恰好等于
F(5)=5，\texttt{F(0)} 的次数等于 F(4)=3）。

\end{addbox}

\begin{defbox}[：名词统一]

课件“函数调用栈 /
递归工作栈”，书上“递归工作栈”；课件“递归表达式（递归体）“＝书上”递归体（递归表达式）“；课件”边界条件（递归出口）“＝书上相同。课件说栈里存”实参”，书上说存“形参”：指的都是这一层函数调用时参数的值（如
\texttt{n=10}）。课件函数名 \texttt{Fib(n)}，书上 \texttt{F(n)}。

\end{defbox}

\hypertarget{ux4e09ux77e5ux8bc6ux56deux987e}{%
\section{三、知识回顾}\label{ux4e09ux77e5ux8bc6ux56deux987e}}

\slidefig{assets/3.3.3-p08.jpg}{0.55}{课件：知识回顾与重要考点}

\begin{summarybox}

\textbf{函数调用的特点}：最后被调用的函数最先执行结束（LIFO）。

\textbf{函数调用栈存储}：① 调用返回地址 ② 实参 ③ 局部变量。

\textbf{递归}：把原问题转换为属性相同、规模较小的问题；需要递归表达式（递归体）和边界条件（递归出口）。递归调用时，函数调用栈可称为\textbf{递归工作栈}：每进入一层递归，就将所需信息压入栈顶；每退出一层，就从栈顶弹出。

\textbf{缺点}：\hlred{效率低}；太多层递归可能导致\hlred{栈溢出}（空间复杂度＝递归深度）；可能包含\hlred{很多重复计算}（如斐波那契）。

\textbf{改造}：可以自定义栈将递归算法改造成非递归算法。

\textbf{自测}：① 函数调用栈里存哪三样东西？② \texttt{factorial(10)}
递归工作栈最多有几层？③ 递归求 F(5)，F(2) 被计算几次？

\end{summarybox}
\end{document}
